\documentclass[journal]{IEEEtai}

\usepackage[colorlinks,urlcolor=blue,linkcolor=blue,citecolor=blue]{hyperref}

\usepackage{amsmath}
\usepackage{amssymb}

\usepackage{color,array}

\usepackage{graphicx}

%% change the following later - leftover from template
%% \jvol{XX}
%% \jnum{XX}
%% \paper{1234567}
%% \pubyear{2020}
%% \publisheddate{xxxx 00, 0000}
%% \currentdate{xxxx 00, 0000}
%% \doiinfo{TQE.2020.Doi Number}


\setcounter{page}{1}
%% \setcounter{secnumdepth}{0}


\begin{document}


\title{Fine-Tuning for Age-Appropriate Pediatric Hospital Communication via LoRA}


\author{Anthony Barros and Kelsey Kloosterman}

\maketitle

\begin{abstract}
Pediatric patients navigating hospital environments require age-appropriate, institutionally grounded information that general-purpose AI systems are not designed to provide. We present an age-targeted fine-tuning framework for a virtual hospital guide deployed within a digital recreation of Victoria Hospital (London, Ontario), specialized for children aged 5--11 and adolescents aged 12--18. Using Low-Rank Adaptation (LoRA) on 4-bit quantized Llama 3.2 Instruct models (1B and 3B) running on an Apple Mac Mini M4 (16~GB) --- a deliberate constraint modeling resource-limited clinical deployment --- we train two adapter configurations: Standard ($r=8$, 16 layers) and Fast ($r=4$, 8 layers), across both model sizes and age groups for 8 training runs total. Training data consists of 500 synthetic examples per adapter generated by Claude Opus 4.6 using Victoria Hospital's public patient materials as grounding artifacts, following a LIMA-inspired curation approach \cite{zhou2023lima}. Fine-tuned adapters substantially outperform zero-shot base models on readability (72--84\% of 5--11 responses meeting a Flesch-Kincaid $\leq 7.0$ grade-level target vs.\ 12--14\% for base models) and inter-role style separation (TF-IDF + logistic regression classifier accuracy 0.89--0.94 vs.\ 0.66--0.70). A configuration-ordering crossover is observed: the Standard adapter yields better results on the 3B model while the Fast adapter performs better on the 1B model, suggesting optimal adapter capacity is model-size-dependent. A subsequent controlled rank sweep (fixed \texttt{num\_layers~=~8}, $r \in \{2, 4, 8, 16\}$, two seeds, four architectures) identifies rank as the operative factor via capacity regularization, finding that transformer depth --- not total parameter count --- predicts the optimal rank regime. Cross-architecture evaluation on Qwen~2 0.5B and SmolLM2 360M confirms the crossover is not Llama-specific; SmolLM2 Standard achieves 84\% FK pass rate and 0.950 classifier accuracy at 0.81\,s average latency despite only 360\,M parameters. These results demonstrate that compact, parameter-efficient adaptation on consumer hardware is feasible for pediatric clinical NLP, and that architecture depth is a more reliable predictor of optimal LoRA rank than total parameter count.
\end{abstract}

\section{Introduction}

\IEEEPARstart{L}{arge} language models (LLMs) have created new opportunities for domain-specific conversational systems, but deploying them in resource-constrained clinical environments introduces practical constraints that are rarely addressed together: limited memory budgets favor smaller models, domain adaptation requires fine-tuning on specialized data, and real-time interaction imposes strict latency bounds.

This paper presents \textit{Dr. Beary Good}, an age-targeted conversational guide embedded in a virtual recreation of Victoria Hospital (London, Ontario). The system provides age-appropriate hospital information to two patient groups --- children aged 5--11 and adolescents aged 12--18 --- via natural language interaction within a VR environment. All training and inference runs on an Apple Mac Mini M4 (16~GB), a deliberate constraint representing deployment on consumer hardware without cloud infrastructure.

To specialize the system for each age group within this memory budget, we fine-tune Llama 3.2 Instruct models \cite{dubey2024llama3} at 1B and 3B scales using Low-Rank Adaptation (LoRA) \cite{hu2022lora}. Training data is generated synthetically following the LIMA hypothesis \cite{zhou2023lima}: 500 training examples per adapter grounded in Victoria Hospital's public patient and family materials, rather than a large noisy corpus. We compare two adapter configurations (Standard: $r = 8$, 16 adapted layers; Fast: $r = 4$, 8 adapted layers) across both model sizes and age groups for eight total training runs.

Our evaluation surfaces a \textbf{configuration-ordering crossover}: the Standard adapter yields better readability and style-separation results on the 3B model, while the Fast adapter outperforms on the 1B model. This crossover is consistent across all five readability metrics and the inter-role classification task, and suggests that optimal adapter capacity is model-size-dependent --- a dependency not captured by common LoRA guidelines that recommend a fixed rank range without model-size conditioning \cite{biderman2024lora}. Only the Fast-1B configuration meets the 1.0-second real-time latency target for VR interaction, specifically for the 5--11 age group (0.93s average); the 12--18 adapter averages 1.20s.

\medskip\noindent The contributions of this paper are:
\begin{itemize}
    \item A complete age-adaptive fine-tuning pipeline for pediatric hospital communication, trained and evaluated on consumer hardware.
    \item Empirical evidence of a LoRA configuration-ordering crossover between Llama 3.2 1B and 3B, consistent across readability and style-separation metrics.
    \item Mechanism identification via a controlled rank sweep (32 runs: 4 ranks $\times$ 4 architectures $\times$ 2 seeds), confirming rank-driven capacity regularization and establishing that transformer depth predicts optimal rank regime.
    \item Cross-architecture validation on Qwen~2 0.5B and SmolLM2 360M, with an updated deployment analysis identifying Qwen Fast and SmolLM2 Standard as strong alternatives to Fast-1B.
\end{itemize}

\section{Related Work}

\subsection{Parameter-Efficient Fine-Tuning}

Low-Rank Adaptation (LoRA) \cite{hu2022lora} reduces fine-tuning cost by constraining weight updates to a low-rank decomposition $\Delta W = BA$, leaving the pre-trained weights frozen. QLoRA \cite{dettmers2023qlora} extended this to 4-bit quantized base models by combining NormalFloat quantization with LoRA, enabling fine-tuning of 7B--65B models on consumer hardware. Both methods are now standard practice in resource-constrained deployment pipelines.

Rank selection in LoRA is typically treated as a task-dependent hyperparameter. Common practitioner guidance --- including the HuggingFace PEFT library defaults and QLoRA \cite{dettmers2023qlora} --- recommends rank 8--16 without conditioning on model scale. Biderman et al.~\cite{biderman2024lora} provide the most thorough empirical study to date, finding that full fine-tuning produces weight perturbations 10--100$\times$ higher in rank than typical LoRA configurations and recommending rank 16 as a practical default across coding and math tasks. This guidance is developed and validated on models at the 7B--70B scale; rank ordering behavior at sub-7B scales is not characterized by existing work. Additionally, the \texttt{scale} (alpha) convention differs across frameworks --- some use the ratio alpha/rank, others a raw scalar --- creating reproducibility risks when configurations are transferred across codebases. Our work provides empirical evidence that adapter performance ordering can reverse between 1B and 3B models on the same task, a regime not covered by prior rank selection studies.

\subsection{LLMs in Clinical and Health Communication}

LLMs have been applied to patient education, clinical FAQ systems, and virtual health assistants. Pediatric deployment introduces challenges distinct from adult-facing systems: responses must simultaneously match developmental vocabulary (shorter sentences, concrete rather than abstract explanations), maintain an emotionally supportive register (reassuring rather than clinical), and enforce hard safety constraints (never implying a diagnosis, always redirecting emergency questions). These properties are not delivered by generic instruction-following fine-tuning on a general-purpose model.

A recurring constraint in clinical NLP is the unavailability of naturalistic patient dialogue due to privacy regulations, motivating synthetic data generation. The LIMA hypothesis \cite{zhou2023lima} --- that a small set of carefully curated, high-quality examples is sufficient for behavioral alignment --- is particularly relevant here: it supports viable training corpora grounded in institutional reference materials without large-scale data collection or patient contact. Our dataset follows this approach, generating 500 examples per adapter from a frontier model conditioned on Victoria Hospital's public patient and family guide. To our knowledge, no prior work has used fine-tuned, age-stratified LoRA adapters for pediatric hospital communication; existing approaches rely on prompted general-purpose models without adapter-level age targeting.

\subsection{Virtual Reality in Pediatric Care}

VR-based interventions have demonstrated measurable efficacy in reducing procedural pain and anxiety in pediatric patients across a range of clinical settings, from preoperative preparation to painful bedside procedures. Despite this evidence base, existing pediatric VR deployments are largely passive: they deliver scripted or pre-rendered content as a distraction mechanism and do not respond to patient queries or adapt to developmental stage. The gap between VR's demonstrated potential and the informational needs of hospitalized children --- who frequently have unanswered questions about their environment, staff roles, and upcoming procedures --- motivates a conversational interface.

Our work addresses this gap by integrating age-targeted fine-tuned adapters into an interactive VR environment. Unlike passive distraction systems, \textit{Dr. Beary Good} responds to patient-initiated queries with responses calibrated to the child's developmental stage, combining the anxiety-reduction potential of immersive VR with the informational utility of a conversational agent grounded in hospital-specific knowledge.

\section{Dataset}
The dataset used to train the models for this project was generated by Claude Opus 4.6 in extended thinking mode. Claude was provided with the Children's Hospital Patient and Family Guide \cite{hospitalguide}, which contains information about the hospital, its routines, procedures, and answers questions that new patients and caretakers may have about their time at the hospital.

From this guide, question and answer pairs were created for five different categories. These categories include Emotional Reassurance, FAQs, Hospital Rules/Routines, What to Expect, and "Who Are These People?". Questions and age-appropriate answers were generated for each of the different age groups, ages 5-11 and 12-18.

For each role and category combination, 100 question and answer pairs were created. For example, the "Emotional Reassurance" category would have 200 questions total, 100 for the 5-11 age group, and 100 for the 12-18 age group. With five different categories, each having entries for both roles, this adds up to 1000 examples total in our dataset.

These 1000 examples are placed in five different .jsonl files, each corresponding to one of the five categories. Each .jsonl file includes all 200 examples for its respective category.

These examples are formatted as follows:

\begin{verbatim}
{"instruction": "...", "response": "...",
"role": "5-11|12-18", "category": "..."}
\end{verbatim}

The "instruction" key holds the user's question, and the "response" key holds the appropriate answer to this question. The "role" key denotes the age group of the user asking the question- this will always be "5-11" or "12-18". Lastly, the "category" key holds the category of the question asked, which would be one of the five categories mentioned prior.

One important thing to note in this dataset is the lack of a system prompt. This is a deliberate choice when desigining both the dataset and the project at large- as this guide will be part of the larger Virtual Hospital project, we must ensure it will fit smoothly with the rest of the project. As the team working on the VR and spatial navigation aspects of the Virtual Hospital already relies on their own system prompt for navigation, refraining from including our own system prompt avoids any potential conflicts.

In addition to this, when using a system prompt, we risk the model simply learning the prompt, and becoming overly reliant on it. When the system prompt is removed, this could lead to a degradation in both quality and the needed tone in the model's responses. By avoiding the use of a system prompt in the first place, we can ensure that the model is learning the needed tone in its weights, and not just following the prompt.

\section{Methodology}\label{sec:methodology}

\subsection{System Overview}

This work extends a virtual hospital framework designed to support safe, naturalistic communication between pediatric patients, caregivers, and clinical decision-makers. The platform instantiates a digital recreation of Victoria Hospital (London, Ontario) in which users interact with an AI-powered guide. The conversational interface routes each query to one of two age-targeted response modules: one serving younger children (ages 5--11) and one serving adolescents (ages 12--18), determined by the established user profile. Emotional-support queries are in scope only to the extent of age-appropriate reassurance and redirection; clinical assessment, therapeutic guidance, and crisis intervention are explicitly out of scope. No system prompt is used: the adapter weights encode the age-appropriate register directly, ensuring that learned behaviour is not prompt-dependent and avoiding conflicts with the VR application's own navigation system prompt.

\subsection{Base Models and Infrastructure}

We use the Llama 3.2 Instruct model family \cite{dubey2024llama3} as our foundation, specifically the 1B and 3B parameter variants, loaded in 4-bit quantized form via the \texttt{mlx\_lm} framework \cite{apple2023mlx} on Apple Silicon. All training and evaluation runs on an Apple Mac Mini M4 (16 GB unified memory) --- a deliberate constraint modeling deployment in resource-limited clinical settings without cloud infrastructure. A lower-bound estimate for full-parameter fine-tuning of the 3B model ($\sim$30 GB in bfloat16 before activations) exceeds the 16 GB ceiling by nearly 2$\times$; the 1B model is borderline with no headroom. Parameter-efficient adaptation via LoRA is therefore the only viable fine-tuning strategy on this hardware for both model scales.

\subsection{LoRA as Declarative Programming}

We employ Low-Rank Adaptation (LoRA) \cite{hu2022lora} to specialize the base models without modifying full model weights. A pre-trained weight matrix $W_0 \in \mathbb{R}^{d \times k}$ is held frozen; a low-rank update is expressed as:

\begin{equation}
\Delta W = BA
\end{equation}

where $B \in \mathbb{R}^{d \times r}$ and $A \in \mathbb{R}^{r \times k}$, with rank $r \ll \min(d, k)$. We insert adapters into all linear projection modules within the targeted transformer layers, using \texttt{scale = 20.0} held constant across all runs. In the \texttt{mlx\_lm} framework, \texttt{scale} is a raw output multiplier applied after the low-rank update --- it is not the alpha/rank ratio used in some other implementations --- so scale~=~20.0 applies identically to both configurations regardless of rank.

We conceptualize LoRA adapters as a form of \textit{declarative programming} over foundation models: hospital guidelines, care pathways, and age-appropriate communication standards collectively define a behavioral specification that the adapter realizes during inference, without encoding knowledge as explicit procedural rules.

\textbf{Adapter configurations evaluated.} Two configurations were trained and compared:

\begin{itemize}
    \item \textbf{Standard adapter} (rank $r = 8$, \texttt{num\_layers = 16}): covers the upper 16 transformer layers --- 57\% of the 3B model and 100\% of the 1B model.
    \item \textbf{Fast adapter} (rank $r = 4$, \texttt{num\_layers = 8}): covers the upper 8 layers --- 29\% of the 3B model and 50\% of the 1B model. Prioritizes efficiency: fewer trainable parameters and faster inference.
\end{itemize}

The two configurations differ in both rank and layer count simultaneously, making it impossible to attribute performance differences to either factor in isolation; this is a known limitation of the comparison design.

\subsection{Synthetic Data Generation}

Following the LIMA hypothesis \cite{zhou2023lima} --- that a small number of carefully curated examples suffices for alignment --- we construct synthetic training corpora rather than collecting naturalistic patient data, motivated both by the impracticality of obtaining real pediatric clinical dialogues at scale and by ethical constraints. Training data was generated by Claude Opus 4.6 (Anthropic; March 2026) with Victoria Hospital's publicly available patient and family reference materials attached as grounding artifacts. Data is organized into five topic categories:

\begin{itemize}
    \item \textbf{What to Expect} --- medical procedures, equipment, and sensory experiences
    \item \textbf{Who Are These People?} --- age-appropriate explanations of clinical staff roles
    \item \textbf{Hospital Rules \& Routines} --- visiting hours, daily schedules, and preparation
    \item \textbf{Emotional Reassurance} --- responses to fear or confusion (non-diagnostic)
    \item \textbf{FAQs / General Curiosity} --- open-ended questions a child or caregiver might ask
\end{itemize}

Each adapter serves a single age group, with 500 training examples per adapter across five categories. Hard safety constraints were enforced at generation time: the model was instructed never to imply a diagnosis, never to recommend medications, and to redirect emergency questions to clinical staff. A held-out validation set of 50 examples per age group was generated in a separate pass. Because both sets share the same generative process, perplexity measures in-distribution coherence rather than out-of-distribution generalization.

\subsection{Training Protocol}

Adapters are trained using the Adam optimizer with a cosine decay learning rate schedule: peak learning rate $1 \times 10^{-5}$, linear warmup over the first 100 steps, decaying to zero over 600 total steps. Batch size is 4; maximum sequence length is 768 tokens; LoRA dropout is 0.0. Loss is computed only over assistant response tokens (\texttt{mask\_prompt: true}). With 500 training examples and batch size 4, 600 steps corresponds to approximately 5 passes over the training data. Both configurations are trained across both model sizes and both age groups, yielding 8 total runs (2 configurations $\times$ 2 model sizes $\times$ 2 age groups), all using seed 42.

\subsection{Evaluation Framework}

Three categories of evaluation metrics are reported. \textbf{Readability:} five formula-based metrics are computed over model-generated responses using the \texttt{textstat} library \cite{bansal2023textstat}. Flesch-Kincaid grade level \cite{kincaid1975fk} is the primary metric; a pass/fail threshold of FK~$\leq$~7.0 is applied to 5--11 adapter outputs, targeting the upper bound of 6th-grade readability while accommodating inherently multi-syllabic clinical vocabulary (e.g., ``Spiritual Care Specialist''). SMOG \cite{mclaughlin1969smog} is designed specifically for health communication text and evaluates polysyllabic word density; it is reliable only on longer responses and returns~0 for very short texts. Gunning Fog \cite{gunning1952fog} penalises words of three or more syllables. Coleman-Liau \cite{coleman1975cli} is character-based rather than syllable-based, providing a formula-independent cross-check. Type-token ratio (TTR) measures lexical diversity (unique words divided by total words; higher values indicate more varied vocabulary) and is reported as a descriptive characteristic. \textbf{Latency and throughput:} per-response latency and tokens-per-second are measured under sequential, single-request conditions on the Mac Mini M4 (no batching, matching VR deployment conditions); a 1.0-second average threshold reflects real-time interaction requirements. Token count is measured from the tokenizer output. \textbf{Inter-role style separation:} a TF-IDF~$+$ logistic regression classifier (scikit-learn \cite{pedregosa2011sklearn}, TF-IDF with max 5{,}000 features and unigram+bigram range, LR with max 1{,}000 iterations) is trained on the combined outputs from both adapters --- each adapter evaluated on its 50 age-matched validation prompts --- using 5-fold cross-validation; accuracy~$\geq$~0.90 is interpreted as strong register separation (chance~=~0.50).

\section{Results}

This section reports results for the eight training runs described in Section 4.5: two adapter configurations (Standard, $r = 8$, \texttt{num\_layers = 16}; Fast, $r = 4$, \texttt{num\_layers = 8}) across two model sizes (1B, 3B) and two age groups (5--11, 12--18). All results correspond to the final checkpoint at step 600. Base model (no-adapter) results are included as a zero-shot reference.

\subsection{Training Dynamics and Perplexity}

Table~\ref{tab:res1} reports post-hoc validation perplexity for all eight Llama adapter configurations. Values are computed via a full-sequence cross-entropy pass over the validation set (50 examples per age group) using the step-600 checkpoint; the loss includes all tokens, unlike the masked training objective. All eight adapters converge without divergence.

The perplexity results mirror the configuration-ordering crossover observed in readability. For the 3B base model, Standard achieves lower perplexity than Fast on both age groups (5--11: 18.17 vs.\ 20.77; 12--18: 16.95 vs.\ 21.61), indicating a better-fit adapter. For the 1B base model, the ordering reverses: Fast achieves lower perplexity than Standard (5--11: 22.32 vs.\ 24.17; 12--18: 19.79 vs.\ 22.31). This corroborates the FK $\leq 7.0$ pass-rate crossover and strengthens the capacity-regularization interpretation: Standard (higher rank) overfits relative to the 1B model's capacity, while Fast (lower rank) underfits relative to the 3B model.

\begin{table}[!t]
  \centering
  \caption{Validation perplexity at step 600 for all Llama adapter configurations (post-hoc full-sequence loss pass).}
  \label{tab:res1}
  \small
  \begin{tabular}{llcc}
    \hline
    Config. & Age & Val Loss & PPL \\
    \hline
    Standard-3B & 5--11  & 2.900 & 18.17 \\
    Standard-3B & 12--18 & 2.830 & 16.95 \\
    Standard-1B & 5--11  & 3.185 & 24.17 \\
    Standard-1B & 12--18 & 3.105 & 22.31 \\
    Fast-3B     & 5--11  & 3.034 & 20.77 \\
    Fast-3B     & 12--18 & 3.073 & 21.61 \\
    Fast-1B     & 5--11  & 3.106 & 22.32 \\
    Fast-1B     & 12--18 & 2.985 & 19.79 \\
    \hline
  \end{tabular}
\end{table}

\subsection{Readability}

Table~\ref{tab:res2} reports readability metrics for all six configurations. All four fine-tuned adapters substantially outperform the base models on the FK $\leq 7.0$ target for the 5--11 group (72--84\% pass rate vs.\ 12--14\% for base models), driven largely by shorter response length (72--99 avg words vs.\ 152--210 for base models). The 12--18 adapters produce consistently higher FK grade levels than the 5--11 adapters (gap: 2.4--3.2 FK points), confirming that the two adapters have learned distinct communication registers.

A \textbf{configuration-ordering crossover} is visible in Table~\ref{tab:res2}: for the Standard configuration, Standard-3B outperforms Standard-1B on FK $\leq 7.0$ pass rate (84\% vs.\ 72\%); for the Fast configuration, this ordering reverses --- Fast-1B outperforms Fast-3B (82\% vs.\ 76\%). This crossover is consistent across all five readability metrics for the 5--11 group and is the central empirical observation of this evaluation. For the four grade-level metrics (FK, SMOG, Gunning Fog, Coleman-Liau), lower values indicate greater age-appropriateness for the 5--11 group; for TTR, higher values indicate more lexically varied responses, which is preferable. The crossover is directionally consistent under both interpretations.

\begin{table*}[!t]
  \caption{Readability metrics (outputs from age-matched held-out prompts, 50 examples per group).}
  \label{tab:res2}
  \small
  \begin{tabular}{llccccccc}
    \hline
    Configuration & Age Group & FK Grade & SMOG & Gunning Fog & Coleman-Liau & TTR & Avg Words & FK $\leq$ 7.0 \\
    \hline
    Standard-1B & 5--11  & 6.1  & 8.8  & 7.7  & 6.7  & 0.712 & 74  & 36/50 (72\%) \\
    Standard-1B & 12--18 & 8.8  & 11.7 & 11.3 & 10.3 & 0.701 & 90  & ---          \\
    Standard-3B & 5--11  & 5.5  & 8.5  & 7.4  & 6.3  & 0.724 & 75  & 42/50 (84\%) \\
    Standard-3B & 12--18 & 8.7  & 11.4 & 10.9 & 10.0 & 0.714 & 92  & ---          \\
    Fast-1B     & 5--11  & 5.5  & 8.4  & 7.2  & 6.2  & 0.701 & 72  & 41/50 (82\%) \\
    Fast-1B     & 12--18 & 8.4  & 11.3 & 10.6 & 9.4  & 0.650 & 92  & ---          \\
    Fast-3B     & 5--11  & 5.9  & 8.9  & 7.9  & 6.6  & 0.636 & 99  & 38/50 (76\%) \\
    Fast-3B     & 12--18 & 8.3  & 11.2 & 10.7 & 9.7  & 0.613 & 122 & ---          \\
    Base-1B     & 5--11  & 9.5  & 12.1 & 12.2 & 9.0  & 0.508 & 201 & 7/50 (14\%)  \\
    Base-1B     & 12--18 & 10.8 & 13.1 & 13.4 & 11.2 & 0.527 & 210 & ---          \\
    Base-3B     & 5--11  & 9.4  & 12.1 & 12.2 & 9.6  & 0.609 & 152 & 6/50 (12\%)  \\
    Base-3B     & 12--18 & 10.6 & 12.9 & 13.3 & 11.4 & 0.565 & 206 & ---          \\
    \hline
  \end{tabular}
\end{table*}

\subsection{Latency and Throughput}

Table~\ref{tab:res3} reports inference latency for all configurations (sequential, no batching, Mac Mini M4). \textbf{Fast-1B is the only configuration meeting the 1.0-second average latency target} for the 5--11 adapter (0.93s avg, 70\% of responses under target). The Fast-1B 12--18 adapter averages 1.20s, missing the threshold. All 3B configurations exceed 2.3 seconds on average. Fast-1B's latency advantage over Base-1B is driven primarily by output length rather than throughput: Fast-1B generates 88 avg tokens at 94.3 tok/s, while Base-1B generates 247 avg tokens at 122.3 tok/s --- the base model is faster per token but produces responses nearly 3$\times$ longer, making fine-tuning the dominant factor in meeting the latency target. Counterintuitively, Fast-3B is slower than Standard-3B despite fewer adapter parameters: Fast-3B generates $\sim$30\% more tokens per response (120/153 vs.\ 92/116 avg tokens), overriding any throughput benefit from the smaller adapter.

\begin{table*}[!t]
  \caption{Inference latency and throughput (50 prompts per age group, sequential, Mac Mini M4).}
  \label{tab:res3}
  \small
  \begin{tabular}{llcccccc}
    \hline
    Config. & Age & Avg Lat. (s) & Min (s) & Max (s) & Avg Tok & Tok/s & $<$1.0s \\
    \hline
    Standard-1B & 5--11  & 1.09 & 0.66 & 1.69 & 89  & 81.7  & 21/50 (42\%) \\
    Standard-1B & 12--18 & 1.36 & 0.62 & 2.02 & 114 & 83.7  & 3/50 (6\%)   \\
    Standard-3B & 5--11  & 2.37 & 1.31 & 3.45 & 92  & 38.8  & 0/50 (0\%)   \\
    Standard-3B & 12--18 & 2.94 & 1.75 & 5.43 & 116 & 39.3  & 0/50 (0\%)   \\
    Fast-1B     & 5--11  & 0.93 & 0.50 & 1.40 & 88  & 94.3  & 35/50 (70\%) \\
    Fast-1B     & 12--18 & 1.20 & 0.75 & 1.79 & 115 & 95.4  & 12/50 (24\%) \\
    Fast-3B     & 5--11  & 2.83 & 1.31 & 6.80 & 120 & 41.9  & 0/50 (0\%)   \\
    Fast-3B     & 12--18 & 3.63 & 0.99 & 7.33 & 153 & 41.7  & 1/50 (2\%)   \\
    Base-1B     & 5--11  & 2.01 & 0.57 & 2.48 & 247 & 122.3 & 2/50 (4\%)   \\
    Base-1B     & 12--18 & 2.14 & 0.25 & 2.46 & 265 & 121.8 & 4/50 (8\%)   \\
    Base-3B     & 5--11  & 3.99 & 0.66 & 6.38 & 190 & 46.5  & 3/50 (6\%)   \\
    Base-3B     & 12--18 & 5.39 & 0.67 & 6.29 & 264 & 48.7  & 1/50 (2\%)   \\
    \hline
  \end{tabular}
\end{table*}

\subsection{Inter-Role Style Separation}

Table~\ref{tab:res4} reports TF-IDF + logistic regression classification accuracy (5-fold CV; chance = 0.50) for distinguishing 5--11 vs.\ 12--18 adapter outputs. All fine-tuned configurations achieve strong separation (0.89--0.94); base models reach only moderate separation (0.66--0.70), confirming that fine-tuning is the primary driver of distinct age registers. The same crossover observed in readability is present here: Standard-3B outperforms Standard-1B (0.920 vs.\ 0.890), while Fast-1B outperforms Fast-3B (0.940 vs.\ 0.900).

\begin{table}[!t]
  \centering
  \caption{Inter-role TF-IDF + LR classification accuracy (5-fold CV; chance = 0.50).}
  \label{tab:res4}
  \small
  \begin{tabular}{lc}
    \hline
    Configuration & Classification Accuracy \\
    \hline
    Standard-1B & $0.890 \pm 0.102$ \\
    Standard-3B & $0.920 \pm 0.051$ \\
    Fast-1B     & $0.940 \pm 0.073$ \\
    Fast-3B     & $0.900 \pm 0.077$ \\
    Base-1B     & $0.660 \pm 0.097$ \\
    Base-3B     & $0.700 \pm 0.063$ \\
    \hline
  \end{tabular}
\end{table}


\subsection{Discussion}

\textbf{Fast-1B is the recommended configuration} for VR deployment: it is the only configuration meeting the 1.0-second latency target (0.93s avg) while maintaining strong readability (FK 5.5 avg) and style separation (0.940 accuracy). The similarity between Standard and Fast on most quality metrics suggests that register adaptation for this task is achievable at rank 4 with 8 adapted layers --- that high-quality, constrained training data may reduce the adapter capacity required to produce measurable style change.

The \textbf{crossover observation} --- Standard adapter better on 3B, Fast adapter better on 1B for both readability and style separation --- is the central empirical finding. We tentatively interpret this as a capacity-regularization effect: the 1B model's limited representational capacity may benefit from the implicit regularization of a smaller adapter, while the 3B model can exploit the additional degrees of freedom of rank 8. However, since rank and layer count co-vary across configurations and all runs use a single seed, this crossover is not statistically confirmed and requires a controlled rank sweep with multiple seeds to establish. An alternative mechanism is over-adaptation of the 1B model under the Standard configuration: the Standard adapter covers 100\% of the 1B model's transformer layers (vs.\ 57\% for 3B), meaning it modifies a proportionally larger share of the model's representational capacity, which may explain the observed ordering reversal independently of rank.

Structural limitations are discussed fully in Section~\ref{sec:limitations}.

\section{Extended Experiments}\label{sec:extended}

Following the evaluation of the original eight runs, two follow-up experiments were conducted to address the structural limitations identified in Section~\ref{sec:limitations}: a controlled rank sweep to isolate the crossover mechanism, and cross-architecture evaluation on two additional model families.

\subsection{Rank Sweep: Mechanism Identification}\label{sec:ranksweep}

The Standard and Fast configurations differ in both \texttt{rank} and \texttt{num\_layers} simultaneously, making it impossible from the original eight runs to attribute the crossover to either factor. To isolate rank, we fixed \texttt{num\_layers~=~8} and swept $r \in \{2, 4, 8, 16\}$ across four architectures (Llama~1B, Llama~3B, Qwen~2 0.5B, SmolLM2~360M) with two seeds (42 and 1337), targeting the \texttt{age\_5\_11} role only. This yields 32 training runs.

Table~\ref{tab:ranksweep} reports FK~$\leq$~7.0 pass rate as a function of rank averaged across seeds. The four models exhibit two distinct profiles:

\begin{itemize}
    \item \textbf{Decreasing (small-model regime):} Llama~1B peaks at $r=2$ (81\%) and degrades monotonically at higher rank. Qwen~0.5B likewise peaks at $r=2$ (78\%).
    \item \textbf{Increasing (large-model regime):} Llama~3B improves through $r=8$ (69\%) before plateauing. SmolLM2~360M improves monotonically through $r=16$ (76\%) and incurs the largest penalty at $r=2$ (49\%).
\end{itemize}

Because \texttt{num\_layers} is held constant across all sweep runs, the crossover between $r=4$ and $r=8$ cannot be attributed to layer coverage. \textbf{Rank is the operative variable}, consistent with a capacity-regularization mechanism: models with limited representational capacity benefit from the implicit constraint of lower rank; deeper models can exploit additional rank for stronger style adaptation.

Notably, SmolLM2~360M falls in the large-model regime despite having fewer total parameters than any Llama configuration. Its 32 transformer layers and hidden dimension of 960 provide per-layer capacity comparable to deeper models, suggesting that \textbf{transformer depth, not total parameter count, determines rank sensitivity}.

\begin{table}[!t]
  \centering
  \caption{Rank sweep FK $\leq 7.0$ pass rate (\%, mean $\pm$ std across two seeds). \texttt{num\_layers~=~8} fixed. Best per model in \textbf{bold}.}
  \label{tab:ranksweep}
  \small
  \begin{tabular}{lcccc}
    \hline
    Model        & $r=2$               & $r=4$     & $r=8$               & $r=16$ \\
    \hline
    Llama 1B     & \textbf{81 $\pm$ 1} & 65 $\pm$ 1 & 63 $\pm$ 5         & 63 $\pm$ 3 \\
    Llama 3B     & 56 $\pm$ 4          & 60 $\pm$ 4 & \textbf{69 $\pm$ 9} & 68 $\pm$ 0 \\
    Qwen 0.5B    & \textbf{78 $\pm$ 0} & 72 $\pm$ 4 & 75 $\pm$ 3         & 68 $\pm$ 2 \\
    SmolLM2 360M & 49 $\pm$ 7          & 69 $\pm$ 5 & 70 $\pm$ 6         & \textbf{76 $\pm$ 8} \\
    \hline
  \end{tabular}
\end{table}

\subsection{Cross-Architecture Validation}\label{sec:crossarch}

To test whether the crossover generalizes beyond Llama, we evaluate Qwen~2 0.5B (4-bit quantized, ChatML template) and SmolLM2 360M \cite{allal2025smollm2} (4-bit quantized, ChatML template) under the Standard ($r=8$, \texttt{num\_layers~=~16}) and Fast ($r=4$, \texttt{num\_layers~=~8}) configurations, both age groups, seed 42. SmolLM2's chat template injects a default system message when no system role is present; training data was generated with empty-content system messages to suppress this, and inference uses a matching empty system context to maintain consistency with training conditions.

Tables~\ref{tab:crossarch} and~\ref{tab:crossarchclf} summarize readability, latency, and classifier results. Several findings follow:

\textbf{Qwen 0.5B: Standard wins on FK, Fast wins on classifier.} This split is resolved by the rank sweep (Table~\ref{tab:ranksweep}): Qwen's monotonically decreasing rank curve places it in the small-model regime. Both configurations clear the 1.0s latency target (100\% and 98\% of age~5--11 responses, respectively) due to high token throughput ($\sim$175 tok/s).

\textbf{SmolLM2 Standard dominates on all metrics.} It achieves 84\% FK pass rate and 0.950 classifier accuracy --- matching Standard-3B on FK and ranking second overall on the classifier --- while averaging 0.81\,s for the age~5--11 role (88\% under target). The Fast adapter, by contrast, produces substantially longer responses (111 vs.\ 76 avg tokens) and falls to 64\% FK pass rate, indicating rank~4 is under-parameterized for its 32-layer architecture. SmolLM2 Standard represents a strong quality-latency trade-off: the depth of a large model with the throughput of a small one.

\textbf{Updated deployment recommendation.} The optimal configuration depends on priority:
\begin{itemize}
    \item \textit{Lowest latency:} Qwen Fast (0.46\,s avg; 98\% under 1.0\,s; 0.960 classifier).
    \item \textit{Guaranteed sub-1s:} Qwen Standard (100\% under 1.0\,s; 74\% FK; 0.940 classifier).
    \item \textit{Best quality-latency balance:} SmolLM2 Standard (84\% FK; 0.950 classifier; 0.81\,s avg).
    \item \textit{Llama-only deployment:} Fast-1B (0.93\,s avg; 82\% FK; 0.940 classifier).
\end{itemize}

\begin{table}[!t]
  \centering
  \caption{Cross-architecture readability and latency for the age~5--11 role (50 examples, seed 42).}
  \label{tab:crossarch}
  \small
  \begin{tabular}{llccccc}
    \hline
    Config. & FK & FK $\leq$7.0 & Avg Lat. & $<$1.0s & Avg Tok \\
    \hline
    Qwen Std.    & 6.2 & 37/50 (74\%) & 0.59\,s & 50/50 (100\%) & 88 \\
    Qwen Fast    & 5.8 & 34/50 (68\%) & 0.46\,s & 49/50 (98\%)  & 83 \\
    SmolLM2 Std. & 5.8 & 42/50 (84\%) & 0.81\,s & 44/50 (88\%)  & 76 \\
    SmolLM2 Fast & 6.1 & 32/50 (64\%) & 1.08\,s & 34/50 (68\%)  & 111 \\
    \hline
  \end{tabular}
\end{table}

\begin{table}[!t]
  \centering
  \caption{Inter-role classifier accuracy for cross-architecture configurations (5-fold CV; chance~=~0.50).}
  \label{tab:crossarchclf}
  \small
  \begin{tabular}{lc}
    \hline
    Configuration    & Accuracy \\
    \hline
    Qwen Standard    & $0.940 \pm 0.049$ \\
    Qwen Fast        & $0.960 \pm 0.058$ \\
    SmolLM2 Standard & $0.950 \pm 0.032$ \\
    SmolLM2 Fast     & $0.920 \pm 0.040$ \\
    \hline
  \end{tabular}
\end{table}

\section{Limitations}\label{sec:limitations}

\textbf{Model family and scale.} The original eight runs use Llama 3.2 at 1B and 3B scales only. The extended experiments (Section~\ref{sec:extended}) add Qwen~2 0.5B and SmolLM2 360M, confirming the crossover is not Llama-specific. Models above 3B, the Qwen 2.5 family, and the Gemma and Mistral families remain untested.

\textbf{Confounded adapter configurations (partially resolved).} The original Standard and Fast configurations differ in both rank and \texttt{num\_layers}. The rank sweep (Section~\ref{sec:ranksweep}) holds \texttt{num\_layers} fixed and confirms rank as the operative factor for Llama and Qwen. The independent effect of \texttt{num\_layers} on the original Llama configurations --- specifically whether Standard-1B's 100\% layer coverage is a contributing factor --- remains unquantified without a dedicated layer sweep.

\textbf{Single-seed runs (partially resolved).} The original eight runs use seed 42. The rank sweep uses two seeds (42 and 1337) and reports mean $\pm$ std; variance estimates for the cross-architecture runs and the original eight configurations are not available.

\textbf{Synthetic data and in-distribution evaluation.} Both training and validation data are generated by the same model under the same prompting protocol, and the 50 validation prompts per age group are drawn from the same topic distribution as training. Perplexity measures in-distribution coherence; performance on novel scenarios or out-of-distribution queries is untested. Human-written reference responses and external evaluators are needed to confirm that adapters produce genuinely age-appropriate output rather than learned stylistic imitation.

\textbf{No clinical validation.} The system has not been tested with pediatric patients, caregivers, or clinical staff. \textit{Dr. Beary Good} is a research prototype; suitability for deployment requires clinical usability studies and institutional review.

\textbf{Safety evaluation.} No automated safety evaluation was performed in this work. A guard model for output filtering is planned for future deployment; the current system relies on the safety alignment of the Llama 3.2 Instruct base model and the safety constraints enforced at data generation time.

\section{Conclusion}

This paper demonstrates that LIMA-style synthetic fine-tuning on consumer hardware can produce measurable, age-appropriate communication behavior in a pediatric hospital context. Fine-tuned adapters substantially outperform zero-shot base models across all measured dimensions: readability pass rates improve from 12--14\% to 72--84\% on the FK $\leq 7.0$ target for the 5--11 group; inter-role style separation improves from 0.66--0.70 to 0.89--0.94 classification accuracy; and the Fast-1B configuration reduces average latency from 2.01s (Base-1B) to 0.93s for the 5--11 group.

The central empirical finding is a \textbf{configuration-ordering crossover}: the Standard adapter ($r = 8$, 16 layers) outperforms the Fast adapter on the 3B model, while the Fast adapter ($r = 4$, 8 layers) outperforms on the 1B model. This crossover is consistent across all five readability metrics and the inter-role classification task. A subsequent controlled rank sweep (Section~\ref{sec:ranksweep}) confirms \textbf{rank as the operative factor} via capacity regularization, and establishes that transformer depth --- not total parameter count --- predicts which regime a model falls into. Cross-architecture evaluation on Qwen~2 0.5B and SmolLM2 360M confirms the crossover is not Llama-specific.

For deployment, the optimal configuration depends on priority. Qwen Fast (0.46\,s avg; 0.960 classifier) offers the lowest latency. Qwen Standard guarantees 100\% of age~5--11 responses under 1.0\,s. SmolLM2 Standard (84\% FK; 0.950 classifier; 0.81\,s avg) offers the best quality-latency balance. Fast-1B (0.93\,s avg; 0.940 classifier) remains the recommended choice for Llama-only deployments.

Remaining limitations include in-distribution evaluation only, the unquantified independent effect of \texttt{num\_layers} on the original Llama configurations, and the absence of clinical validation with pediatric patients or caregivers (see Section~\ref{sec:limitations}).

\begin{thebibliography}{99}

\bibitem{zhou2023lima}
C.~Zhou, P.~Liu, P.~Xu, S.~Iyer, J.~Sun, R.~Maddela, X.~Peng, S.~Shrivastava, M.~Lewis, L.~Zettlemoyer, and O.~Levy, ``LIMA: Less Is More for Alignment,'' in \textit{Advances in Neural Information Processing Systems (NeurIPS)}, 2023.

\bibitem{dubey2024llama3}
A.~Dubey et al., ``The Llama 3 Herd of Models,'' \textit{arXiv preprint arXiv:2407.21783}, Meta AI, 2024.

\bibitem{apple2023mlx}
Apple, ``MLX: Efficient and Flexible Machine Learning on Apple Silicon,'' 2023. [Online]. Available: https://github.com/ml-explore/mlx

\bibitem{hu2022lora}
E.~J.~Hu, Y.~Shen, P.~Wallis, Z.~Allen-Zhu, Y.~Li, S.~Wang, L.~Wang, and W.~Chen, ``LoRA: Low-Rank Adaptation of Large Language Models,'' in \textit{Proc. Int. Conf. Learning Representations (ICLR)}, 2022.

\bibitem{kincaid1975fk}
J.~P.~Kincaid, R.~P.~Fishburne, R.~L.~Rogers, and B.~S.~Chissom, ``Derivation of New Readability Formulas for Navy Enlisted Personnel,'' Research Branch Report 8-75, Naval Technical Training Command, 1975.

\bibitem{mclaughlin1969smog}
G.~H.~McLaughlin, ``SMOG Grading --- A New Readability Formula,'' \textit{J. Reading}, vol.~12, no.~8, pp.~639--646, 1969.

\bibitem{gunning1952fog}
R.~Gunning, \textit{The Technique of Clear Writing}. New York, NY, USA: McGraw-Hill, 1952.

\bibitem{coleman1975cli}
M.~Coleman and T.~L.~Liau, ``A Computer Readability Formula Designed for Machine Scoring,'' \textit{J. Applied Psychology}, vol.~60, no.~2, pp.~283--284, 1975.

\bibitem{bansal2023textstat}
S.~Bansal, ``textstat: Python Library to Calculate Readability Statistics of a Text Object'' (v0.7.13), 2023. [Online]. Available: \url{https://pypi.org/project/textstat/}

\bibitem{dettmers2023qlora}
T.~Dettmers, A.~Pagnoni, A.~Holtzman, and L.~Zettlemoyer, ``QLoRA: Efficient Finetuning of Quantized LLMs,'' in \textit{Advances in Neural Information Processing Systems (NeurIPS)}, 2023.

\bibitem{biderman2024lora}
D.~Biderman, J.~G.~Ortiz, J.~Portes, M.~Paul, P.~Greengard, C.~Jennings, D.~King, S.~Havens, V.~Chiley, J.~Frankle, C.~Blakeney, and J.~P.~Cunningham, ``LoRA Learns Less and Forgets Less,'' \textit{Transactions on Machine Learning Research (TMLR)}, 2024.

\bibitem{pedregosa2011sklearn}
F.~Pedregosa et al., ``Scikit-learn: Machine Learning in Python,'' \textit{J. Machine Learning Research}, vol.~12, pp.~2825--2830, 2011.

\bibitem{allal2025smollm2}
L.~B. Allal et al., ``SmolLM2: When Smol Goes Big,'' \textit{arXiv preprint arXiv:2502.05737}, Hugging Face, 2025.

\bibitem{hospitalguide}
“Children's Hospital Patient and Family Guide,” Children's Hospital London Health Sciences Centre. https://www.lhsc.on.ca/childrens-hospital/childrens-hospital-patient-and-family-guide (accessed Mar. 26, 2026).

\end{thebibliography}

\clearpage
\section*{Appendix}

\subsection{Repository Link}
\noindent \url{https://github.com/abarros6/small-bear}

\subsection{Group Member Contributions}
\textbf{Anthony Barros}
\begin{itemize}
    \item \textbf{Abstract}
    \item \textbf{Methodology} 
    \item \textbf{Results}
    \item \textbf{Limitations}
    \item \textbf{Proofreading and Citations}
\end{itemize}

\textbf{Kelsey Kloosterman}
\begin{itemize}
    \item \textbf{Introduction}
    \item \textbf{Related Works} 
    \item \textbf{Dataset} 
    \item \textbf{Conclusion}
    \item \textbf{Proofreading and Citations}
\end{itemize}


\subsection{Claude Transcripts}

All synthetic data was generated using Claude Opus 4.6 with extended thinking mode via Claude.ai Project Spaces (March 2026). Two separate project spaces were created: one for training data and one for validation data. Each project space had Victoria Hospital's publicly available patient and family guide attached as a grounding artifact. Within each space, five dedicated chats were run --- one per category --- to maintain context isolation across topics. Training scripts and evaluation code were developed using Claude Code (claude-sonnet-4-6, v2.1.74).

\subsubsection*{Training Data Generation}

\noindent\textbf{System prompt (summarized):} Generate high-quality synthetic training data for two pediatric-facing adapter models (ages 5--11 and 12--18) for deployment in a virtual recreation of Victoria Hospital in London, Ontario. Following the LIMA hypothesis, produce a small set of carefully curated, high-quality examples rather than a large noisy corpus --- 100 examples per category per adapter, 500 total per adapter. Ground all responses in the attached Victoria Hospital reference artifacts.

\noindent\textbf{Output format:}
\begin{verbatim}
{"instruction": "...", "response": "...",
 "role": "5-11|12-18", "category": "..."}
\end{verbatim}

\noindent\textbf{Hard safety rules enforced at generation time:} never imply a diagnosis; never recommend medications or treatments; never use language inappropriate for the target age range; always redirect medical or emergency questions to a healthcare professional or parent/guardian.

\noindent\textbf{Per-category prompts (training):}
\begin{itemize}
  \item \textit{emotional\_reassurance}: ``Generate 200 training examples for the emotional\_reassurance category. Ensure every example that involves a situation requiring medical attention redirects to a nurse, doctor, or parent/guardian. Produce 100 examples tagged `role': `5-11' and 100 tagged `role': `12-18'.''
  \item \textit{faqs\_general\_curiosity}: ``Generate 200 training examples for the faqs\_general\_curiosity category. Use the attached Victoria Hospital reference artifacts where relevant. Produce 100 per age group.''
  \item \textit{hospital\_rules\_and\_routines}: ``Generate 200 training examples for the hospital\_rules\_and\_routines category. Prioritize the attached reference artifacts as the source of truth for policies, schedules, and routines. Produce 100 per age group.''
  \item \textit{who\_are\_these\_people}: ``Generate 200 training examples for the who\_are\_these\_people category. Use the attached reference artifacts to ground roles and descriptions in this specific facility. Produce 100 per age group.''
  \item \textit{what\_to\_expect}: ``Generate 200 training examples for the what\_to\_expect category. Use the attached reference artifacts to ground responses in this facility. Produce 100 per age group.''
\end{itemize}

\subsubsection*{Validation Data Generation}

\noindent\textbf{System prompt (summarized):} Generate high-quality synthetic validation data independent from the training set. The validation set represents 10\% of training data --- 10 examples per age group per category, 100 examples total. Validation examples must not reuse or closely mirror training examples. Every example must meet a gold standard of quality; a weak validation example is worse than no example.

\noindent\textbf{Per-category prompts (validation):}
\begin{itemize}
  \item Each of the five categories: ``Generate 20 validation examples for the \textit{[category]} category. Produce 10 examples tagged `role': `5-11' and 10 tagged `role': `12-18'. Output each example as a JSON object on its own line.''
\end{itemize}

\end{document}
